\noindent\emph{QYBE, unitarity, shift, hexagon.}
The quantum Yang--Baxter equation on
$V_{\mathrm{EK}}^{\otimes 3}$ reduces on $V^{\otimes 3}$
to~\eqref{eq:yangian-qybe}, verified in
Example~\ref{ex:yangian}; the extension to
$V_{\mathrm{EK}}$ follows from the universal property of
$\mathcal{R}(z)$.
Unitarity $S_{12}(z) \cdot (\sigma \circ S_{12}(-z)
\circ \sigma) = \id$ follows from Drinfeld--Jimbo unitarity.
The shift condition
\begin{equation}\label{eq:ek-shift-cond}
  [D \otimes \id + \id \otimes D,\; S(z)]
  \;=\; \partial_z\, S(z)
\end{equation}
follows from the translation-covariance of the Drinfeld
coproduct. The hexagon follows from the pentagon relation
for $\Phi$ (from the Stasheff boundary of
$\overline{\FM}_4^{\mathrm{ord}}(\CC) = K_4$) combined
with the QYBE for $S$, by the standard argument of
Drinfeld~\cite{EK96} lifted to the vertex-algebraic setting.

\medskip
\noindent
\textbf{Part B: the ordered chiral center of the EK quantum
VOA.}
We compute the ordered chiral Hochschild cohomology
$\HH^{*,\mathrm{ord}}_{\mathrm{ch}}
(V_{\mathrm{EK}},\,V_{\mathrm{EK}})$
at low degrees on the formal disk $D$, paralleling the
computation of Example~\ref{ex:yangian} (Part B) but now
using the \emph{full} vertex $R$-matrix $S(z)$ rather than
the Yang $R$-matrix on $V \otimes V$ alone.

\smallskip
\noindent\emph{Degree~$2$: the KZ connection with the full
vertex $R$-matrix.}
At degree~$2$, the KZ connection on
$\Conf_2^{\mathrm{ord}}(\CC)$ with coefficients in
$\End(V_{\mathrm{EK}}^{\otimes 2})$ is
\begin{equation}\label{eq:ek-kz-full}
  \nabla_{\mathrm{KZ}}^S
  = d + \frac{\hbar\,\Omega}{z}\,dz
  + \frac{\hbar^2}{z^2}\,
  \Bigl(\Omega^{(2)} + \tfrac{1}{2}\,\Omega^2\Bigr)\,dz
  + O(z^{-3}),
\end{equation}
where $z = z_1 - z_2$. The connection~\eqref{eq:ek-kz-full}
is obtained from $S(z)$ via the general formula
$\nabla = d + (\partial_z \log S(z))\,dz$ (logarithmic
derivative of the vertex $R$-matrix).

For the classical $V_k(\mathfrak{sl}_2)$ (trivial braiding
$S(z) = \id$), the flat sections are
$\Phi(z) = z^{\Omega/(k+2)}$ with algebraic monodromy.
For the EK quantum VOA ($S(z) \neq \id$), the flat sections
acquire quantum corrections:
\begin{equation}\label{eq:ek-flat-sections}
  \Phi^S(z)
  = z^{\hbar\,\Omega}\,
  \Bigl(I + \frac{\hbar^2\,\Omega^{(2)}}{z}
  + O(z^{-2})\Bigr).
\end{equation}

\emph{The kernel $\ker(\av)$.}
The ordered chiral centre computation parallels
Example~\ref{ex:yangian}; the result is identical on
cohomology. On $V \otimes V$:
$\ker(\av_2)\big|_{V \otimes V}
\cong \bigwedge^2 V = \CC$,
the same one-dimensional space as
equation~\eqref{eq:yangian-ker-av2}.
The quantum corrections from $S(z)$ shift the subleading
terms of the flat sections but do not alter the
$\ker(\av_2)$ dimension. The scalar shadow surviving
$\av_2$ is $\kappa = 3(k+2)/4$ (equation~\eqref{eq:sl2-kappa}),
identical to the classical value. At degree~$3$, the kernel
$\ker(\av_3) \cong \Lambda^3(\fg) = \CC$ carries the
quantum-corrected Drinfeld associator
\begin{equation}\label{eq:ek-ker-av3}
  \Phi_{\mathrm{KZ}}^S - 1
  \;=\; \hbar^2\,\zeta(2)\,[\Omega_{12}, \Omega_{23}]
  + O(\hbar^3),
\end{equation}
whose leading term agrees with the classical KZ associator.

\medskip
\noindent
\textbf{Part C: $S$-locality as an $\Ainf$ relation.}
The braided VOA axioms of the EK quantum VOA are equivalent
to the chiral $\Ainf$ relations for $V_{\mathrm{EK}}$
through the chiral quantum group equivalence
(Theorem~\ref{thm:chiral-qg-equiv},
(I)$\Leftrightarrow$(II)).
In the ordered bar complex
$\Barord(V_{\mathrm{EK}}) = T^c(s^{-1}\overline{V}_{\mathrm{EK}})$,
the $S$-locality axiom~\eqref{eq:ek-s-locality} governs the
reordering of two insertions: the braid group
$\pi_1(\Conf_2^{\mathrm{ord}}(\CC)) = \ZZ$ acts on degree-$2$
bar chains by
$\sigma_1 \cdot [a | b] = S(z) \cdot [b | a]$.
The $S$-locality axiom is the $n = 2$ Stasheff relation,
the hexagon is the $n = 3$ relation, and the QYBE is the
$n = 4$ relation (from the pentagon for $\Phi$ on
$K_4 = \overline{\FM}_4^{\mathrm{ord}}(\RR)$).

\medskip
\noindent
\textbf{Summary.}
\noindent
The EK quantum VOA $(V_{\mathrm{EK}},\, Y,\, S(z))$
provides the first concrete instance of all three structures
of the chiral quantum group equivalence
(Theorem~\ref{thm:chiral-qg-equiv}) exhibited
simultaneously on a single algebra with all axioms verified.
The ordered chiral center carries the full quantum $R$-matrix
$S(z)$ (not just the classical $r$-matrix
$r(z) = \hbar\,\Omega/z$), and the kernel
$\ker(\av_2) \neq 0$ demonstrates that the ordered theory
is strictly richer than the symmetric Beilinson--Drinfeld
chiral homology even at the quantum level.
\end{example}

\begin{example}[Heisenberg: trivial case]
\label{ex:heis-coproduct}
For the Heisenberg algebra $\cH_k$ (class $G$), all three
structures are trivial:
$S(z) = \id$ (the braiding is symmetric);
$m_k^{\mathrm{ch}} = 0$ for $k \geq 3$ (the algebra is
strictly associative in the chiral sense);
$\Delta^{\mathrm{ch}}(J) = J \otimes 1 + 1 \otimes J$ is
cocommutative and strictly coassociative ($\Phi$ acts
trivially). The ordered and symmetric data coincide.
\end{example}

\begin{theorem}[$\cW_{1+\infty}[\Psi]$ as a chiral quantum
group]
\label{thm:w-infty-chiral-qg}
The algebra $\cW_{1+\infty}[\Psi]$ carries a chiral quantum
group datum in the sense of
Theorem~\textup{\ref{thm:chiral-qg-equiv}}: an $R$-matrix,
a chiral coproduct, and an $\Ainf$ structure, all explicit
and originating from the cohomological Hall algebra of the
Jordan quiver. The precise content is as follows.
\begin{enumerate}[label=\textup{(\Roman*)}]
\item \textup{($R$-matrix.)}
  The Maulik--Okounkov $R$-matrix~\cite{MO19} of quantum
  toroidal $\mathfrak{gl}_1$, restricted to the CoHA, is
  \emph{scalar}:
  \begin{equation}\label{eq:gl1-scalar-r}
    R(z) = g(z)
    = \frac{(z - h_1)(z - h_2)(z - h_3)}
    {(z + h_1)(z + h_2)(z + h_3)},
    \qquad h_1 + h_2 + h_3 = 0,
  \end{equation}
  where $h_1, h_2, h_3$ are the equivariant parameters of
  $\CC^3$ subject to the Calabi--Yau constraint. The
  classical limit is
  % AP126: r-matrix with level prefix Psi; Psi=0 -> r=0. Verified.
  $r(z) = \Psi/z$ \textup{(}class~$G$, level prefix
  $\Psi$\textup{)}.
\item \textup{(Chiral coproduct at all spins.)}
  The Drinfeld coproduct on $Y(\widehat{\mathfrak{gl}}_1)$
  gives an explicit spectral coproduct
  $\Delta_z \colon \cW_{1+\infty}[\Psi] \to
  (\cW_{1+\infty}[\Psi] \mathbin{\hat{\otimes}}
  \cW_{1+\infty}[\Psi])((z))$
  satisfying strict coassociativity,
  the spectral counit, and OPE compatibility at all spins.
  The transfer matrix
  $T(u) = 1 + \sum_{n \geq 1} \psi_n\, u^{-n}$
  is scalar \textup{(}$\mathfrak{gl}_1$ has
  rank~$1$\textup{)}, and
  \begin{equation}\label{eq:gl1-drinfeld-coprod}
    \Delta_z(T(u))
    = T(u) \otimes T(u - z).
  \end{equation}
  The general spin-$n$ formula is
  \begin{equation}\label{eq:gl1-coprod-general}
    \Delta_z(\psi_n)
    = \psi_n \otimes 1
    + \sum_{k=0}^{n-1}\; \sum_{m=1}^{n-k}
    \binom{n-k-1}{m-1}\,
    z^{n-k-m}\,
    \psi_k \otimes \psi_m,
  \end{equation}
  where $\psi_0 = 1$.
\item \textup{($\Ainf$ structure.)}
  The chiral $\Ainf$-operations $m_k^{\mathrm{ch}}$ are
  extracted from the ordered bar complex
  $\Barord(\cW_{1+\infty})
  = T^c(s^{-1}\overline{\cW}_{1+\infty})$ via the functor
  $\alpha$ of Theorem~\textup{\ref{thm:chiral-qg-equiv}},
  with input the $R$-matrix~\eqref{eq:gl1-scalar-r}.
\end{enumerate}
The three structures determine each other by
Theorem~\textup{\ref{thm:chiral-qg-equiv}}. The OPE
compatibility axiom~\eqref{eq:ope-compat} is satisfied at
all spins simultaneously by two independent arguments:
the coderivation property on the Koszul-locus bar complex,
and the JKL vertex bialgebra theorem on the CoHA.
\end{theorem}

\begin{proof}
The argument proceeds in seven steps.

\medskip
\noindent
\textbf{Step~1}
(CoHA identification;
Schiffmann--Vasserot~\cite{SV13},
Schiffmann--Mellit--Minets--Vasserot~\cite{SMMV23}).
The critical CoHA of the Jordan quiver is
\[
  \cH_{\mathrm{Jor}}
  = \bigoplus_{n \geq 0}
  H^{GL_n}_{\mathrm{BM},*}
  (\cM_n^{\mathrm{nil}},\, \varphi_W),
  \qquad
  W = \mathrm{Tr}(x[y,z]),
\]
where $\cM_n^{\mathrm{nil}}$ is the nilpotent locus in
the representation variety of the Jordan quiver and
$\varphi_W$ denotes the vanishing-cycle twist.
Schiffmann--Vasserot~\cite{SV13} proved
$\cH_{\mathrm{Jor}} \cong Y^+(\widehat{\mathfrak{gl}}_1)$
by matching generators and relations; the critical
refinement is~\cite{SMMV23}.
Kontsevich--Soibelman~\cite{KS11} established the general
CoHA framework (see also Joyce~\cite{Joyce18}).

\medskip
\noindent
\textbf{Step~2}
(Bialgebra and Drinfeld double; Yang--Zhao~\cite{YZ18a}).
The CoHA carries a bialgebra coproduct
$\Delta_{\mathrm{CoHA}} \colon
\cH_{\mathrm{Jor}} \to
\cH_{\mathrm{Jor}} \otimes \cH_{\mathrm{Jor}}$
from restriction to Levi subalgebras
$GL_{n_1} \times GL_{n_2} \hookrightarrow GL_n$.
Yang--Zhao proved that the Drinfeld double is the full
affine Yangian: $D(\cH_{\mathrm{Jor}}) \cong
Y(\widehat{\mathfrak{gl}}_1)$.

\medskip
\noindent
\textbf{Step~3}
(Yangian--$\cW$-algebra identification;
Proch\'azka--Rap\v{c}\'ak~\cite{PR19}).
The quantum Miura transform gives an isomorphism of
filtered associative algebras
\[
  Y(\widehat{\mathfrak{gl}}_1)
  \;\cong\;
  U(\cW_{1+\infty}[\Psi]),
\]
intertwining the Yangian filtration with the conformal
filtration. The $\psi$-generators map to
$\cW_{1+\infty}$ fields via a triangular change
of basis: $\psi_1 = J$,
$\psi_n = W_n + P_n(J, W_2, \ldots, W_{n-1})$
for $n \geq 2$, where $W_2 = T$ (stress tensor),
$W_3 = W$ (spin-$3$ current), and $P_n$ is a
nonlinear polynomial in lower-spin fields with
coefficients depending on~$\Psi$.

\medskip
\noindent
\textbf{Step~4}
(Drinfeld coproduct at all spins).
The transfer matrix
$T(u) = 1 + \sum_{n \geq 1} \psi_n\, u^{-n}$
is a scalar formal power series in $u^{-1}$
($\mathfrak{gl}_1$ has rank~$1$). The Drinfeld
coproduct~\eqref{eq:gl1-drinfeld-coprod} is the product
of two such series. The
expansion~\eqref{eq:gl1-coprod-general} follows from
$(u - z)^{-m} = u^{-m} \sum_{j \geq 0}
\binom{m+j-1}{j} (z/u)^j$ applied to
$T(u) \otimes T(u-z)$ and collecting the coefficient
of $u^{-n}$
(see~\eqref{eq:yangian-drinfeld-coprod} and
Example~\ref{ex:yangian-coproduct}).
Through spin~$3$:
\begin{align}
  \Delta_z(\psi_1)
  &= \psi_1 \otimes 1
  + 1 \otimes \psi_1,
  \label{eq:coprod-psi1} \\[4pt]
  \Delta_z(\psi_2)
  &= \psi_2 \otimes 1
  + 1 \otimes \psi_2
  + \psi_1 \otimes \psi_1
  + z\,(1 \otimes \psi_1),
  \label{eq:coprod-psi2} \\[4pt]
  \Delta_z(\psi_3)
  &= \psi_3 \otimes 1
  + 1 \otimes \psi_3
  + \psi_1 \otimes \psi_2
  + \psi_2 \otimes \psi_1
  \notag \\
  &\quad
  + z\,(\psi_1 \otimes \psi_1
  + 2 \cdot 1 \otimes \psi_2)
  + z^2\,(1 \otimes \psi_1).
  \label{eq:coprod-psi3}
\end{align}
In the $\cW_{1+\infty}$ field basis (via the Miura
substitution $\psi_1 = J$,
$\psi_2 = T + J^2/(2\Psi)$ of step~3):
\begin{align}
  \Delta_z(J)
  &= J \otimes 1
  + 1 \otimes J,
  \label{eq:coprod-J} \\[4pt]
  \Delta_z(T)
  &= T \otimes 1
  + 1 \otimes T
  + \tfrac{\Psi - 1}{\Psi}\, J \otimes J
  + z\,(1 \otimes J).
  \label{eq:coprod-T}
\end{align}
The coefficient $(\Psi - 1)/\Psi$ in~\eqref{eq:coprod-T}
arises from the Miura inversion $T = \psi_2 - J^2/(2\Psi)$:
the $\psi_1 \otimes \psi_1$ cross-term in
$\Delta_z(\psi_2)$ has coefficient~$1$, and
subtracting $(1/(2\Psi))\,\Delta_z(J)^2
= (J^2/(2\Psi)) \otimes 1 + (1/\Psi)\,J \otimes J
+ 1 \otimes J^2/(2\Psi)$
replaces the coefficient by $1 - 1/\Psi = (\Psi - 1)/\Psi$.
At $\Psi = 1$ (free boson) the cross-term vanishes;
in the classical limit $\Psi \to \infty$ it tends to~$1$.
The Heisenberg OPE
$J(z)\,J(w) \sim \Psi/(z-w)^2$ gives
% AP126: r-matrix with level prefix Psi; Psi=0 -> r=0. Verified.
$r(z) = \Psi/z$ (class~$G$, level prefix $\Psi$).
The general formula~\eqref{eq:gl1-coprod-general} gives
$\Delta_z(\psi_n)$ at every spin in closed form; the
coproduct on the $\cW_{1+\infty}$ field basis
$\{J, T, W, W_4, \ldots\}$ is obtained by inverting
the Miura transform, and the computation is triangular
and algorithmic at each spin.

\medskip
\noindent
\textbf{Step~5}
($R$-matrix from Maulik--Okounkov; vertex bialgebra
structure; Maulik--Okounkov~\cite{MO19},
Jindal--Kaubrys--Latyntsev~\cite{JKL26}).
The critical CoHA of any quiver with potential carries
a vertex coproduct forming a vertex
bialgebra (\cite{JKL26},~Theorems~A and~C). The
$R$-matrix is
\[
  R(z)
  = \frac{\Psi(\sigma^* \mathrm{Ext}^\vee, z)}
  {\Psi(\mathrm{Ext}, z)},
\]
where $\Psi$ on the right denotes the motivic exponential
and $\mathrm{Ext}$ is the $\mathrm{Ext}^1$-bundle on the
stack of quiver representations.
The Maulik--Okounkov $R$-matrix~\cite{MO19}
restricted to the CoHA gives
$R_{\mathrm{CoHA}} = q^{\alpha \cdot \beta}$
($q = e^{\pi i \Psi}$) with classical limit
$r(z) = \Psi/z$.

\medskip
\noindent
\textbf{Step~6}
(OPE compatibility at all spins: two independent
arguments).
The vertex bialgebra axiom~\eqref{eq:ope-compat}
requires $\Delta_z(Y(a,w)\,b)
= Y^{(2)}(\Delta_z(a),w)\,\Delta_z(b)$
as an identity of vertex operators, not merely of
associative algebra elements. We give two independent
proofs that this holds at all spins simultaneously.

\smallskip
\noindent\emph{Argument~A \textup{(}coderivation on the
Koszul locus\textup{)}.}
The algebra $\cW_{1+\infty}[\Psi]$ at generic~$\Psi$ lies
on the Koszul locus
\textup{(}Definition~\textup{\ref{def:koszul-locus}}:
the bar-cobar adjunction
$\Omega(\Barord(\cW_{1+\infty}[\Psi]))
\xrightarrow{\;\sim\;} \cW_{1+\infty}[\Psi]$ is an
equivalence\textup{)}. By the direction
$\textup{(II)} \to \textup{(III)}$ of
Theorem~\textup{\ref{thm:chiral-qg-equiv}}, the bar
construction produces a chiral coproduct on
$\cW_{1+\infty}[\Psi]$ from the deconcatenation coproduct
$\Delta$ on $\Barord(\cW_{1+\infty}[\Psi])
= T^c(s^{-1}\overline{\cW}_{1+\infty}[\Psi])$
via the cobar dualization~\eqref{eq:coprod-from-bar}.
The bar differential $d_{\Barord}$ is a coderivation with
respect to~$\Delta$:
\begin{equation}\label{eq:coderivation-w-infty}
  d_{\Barord} \circ \Delta
  = \Delta \circ d_{\Barord}.
\end{equation}
This identity holds at \emph{all arities
simultaneously} as a structural property of the cofree
tensor coalgebra. Under the cobar functor~$\Omega$, the
coderivation property~\eqref{eq:coderivation-w-infty}
dualizes to the OPE compatibility
axiom~\eqref{eq:ope-compat}: the bar differential
encodes the vertex operation $Y(a,w)$, and the
deconcatenation encodes the coproduct $\Delta_z$
\textup{(}see the proof of
Theorem~\textup{\ref{thm:chiral-qg-equiv}}, direction
$\textup{(II)} \to \textup{(III)}$\textup{)}.

It remains to identify this bar-cobar coproduct with the
Drinfeld coproduct of step~4. The bar-cobar coproduct
is determined by the vertex algebra structure of
$\cW_{1+\infty}[\Psi]$; the Drinfeld coproduct is
determined by the generating-series
formula~\eqref{eq:gl1-drinfeld-coprod}. Both satisfy
coassociativity and the counit axiom. On the Koszul locus,
the coproduct compatible with a given vertex $R$-matrix is
unique (Theorem~\ref{thm:chiral-qg-equiv}: the functors
$\alpha$, $\beta$, $\gamma$ form an equivalence triangle).
The bar-cobar coproduct is compatible with the
Maulik--Okounkov $R$-matrix (by construction via~$\alpha$
and~$\gamma^{-1}$). The Drinfeld coproduct is compatible
with the same $R$-matrix (the Yangian $R$-matrix of
step~5 restricts to the MO $R$-matrix). By uniqueness,
the two coproducts coincide. Therefore the Drinfeld
coproduct satisfies OPE compatibility at all spins.

\smallskip
\noindent\emph{Argument~B \textup{(}JKL vertex bialgebra
on the CoHA\textup{)}.}
Jindal--Kaubrys--Latyntsev~\cite{JKL26} construct a vertex
coproduct on the critical CoHA of any quiver with potential,
satisfying the vertex bialgebra
axiom~\eqref{eq:ope-compat} (their Theorems~A and~C).
For the Jordan quiver, the CoHA is
$\cH_{\mathrm{Jor}} \cong Y^+(\widehat{\mathfrak{gl}}_1)$
(step~1), and the JKL vertex coproduct recovers the
Drinfeld coproduct on the Yangian (their Theorem~B).
Since the JKL vertex bialgebra axiom is exactly the OPE
compatibility axiom~\eqref{eq:ope-compat}, this gives a
second proof at all spins simultaneously.

\smallskip
\noindent\emph{RTT triviality \textup{(}corroboration,
not proof\textup{)}.}
For $\mathfrak{gl}_1$, the Maulik--Okounkov $R$-matrix is
scalar:
$R(z) = g(z)$ of~\eqref{eq:gl1-scalar-r}. The RTT
relation
\begin{equation}\label{eq:rtt-gl1}
  R(u - v)\, T_1(u)\, T_2(v)
  = T_2(v)\, T_1(u)\, R(u - v)
\end{equation}
is trivially satisfied because a scalar commutes through
all operators. This confirms that the Drinfeld coproduct is
compatible with the $R$-matrix at the \emph{associative
algebra} level; arguments~A and~B upgrade this to the
full vertex algebra compatibility.
For $Y(\widehat{\mathfrak{gl}}_N)$ with $N \geq 2$,
the $R$-matrix is no longer scalar, and OPE compatibility
at each truncation is a separate question
(Remark~\ref{rem:w-infty-vertex-gap}).

\medskip
\noindent
\textbf{Step~7}
(Full axiom verification).
The spectral coproduct~\eqref{eq:gl1-drinfeld-coprod}
satisfies all axioms of
Definition~\ref{def:chiral-coproduct}:
\begin{enumerate}[label=\textup{(\alph*)}]
\item \emph{Strict coassociativity.}
  $(\Delta_{z_1} \otimes \id) \circ \Delta_{z_2}(T(u))
  = T(u) \otimes T(u - z_2) \otimes T(u - z_1 - z_2)$
  and
  $(\id \otimes \Delta_{z_2}) \circ \Delta_{z_1}(T(u))
  = T(u) \otimes T(u - z_1) \otimes T(u - z_1 - z_2)$.
  For $\mathfrak{gl}_1$ the Drinfeld associator
  acts trivially (the algebra is abelian), so
  $\Delta_z$ is \emph{strictly} coassociative.
\item \emph{Counit.}
  $\varepsilon(\psi_n) = 0$ for $n \geq 1$,
  $\varepsilon(1) = 1$ gives
  $(\varepsilon \otimes \id) \circ \Delta_z(T(u))
  = T(u - z)$ and
  $(\id \otimes \varepsilon) \circ \Delta_z(T(u))
  = T(u)$.
\item \emph{OPE compatibility at all spins.}
  By step~6: the coderivation property on the bar complex
  of $\cW_{1+\infty}[\Psi]$ (argument~A) gives the vertex
  bialgebra axiom~\eqref{eq:ope-compat} at all spins via
  the Koszul-locus equivalence
  (Theorem~\ref{thm:chiral-qg-equiv}), and the JKL vertex
  bialgebra theorem on the CoHA (argument~B) provides an
  independent proof. RTT triviality for $\mathfrak{gl}_1$
  corroborates the result at the associative algebra level.
\end{enumerate}
No constant coproduct exists for $c \neq 0$: the
conformal anomaly $T(z)\,T(w) \sim (c/2)/(z-w)^4 +
\cdots$ obstructs any $z$-independent ansatz
$\Delta(T) = T \otimes 1 + 1 \otimes T + X$
because the quartic pole on the left maps to
$(c/2 + c/2)/(z-w)^4$ while the right retains $c/2$.
The spectral parameter in~\eqref{eq:gl1-drinfeld-coprod}
absorbs this obstruction through the shift
$u \mapsto u - z$.

The three structures (I)--(III) determine each other by
Theorem~\ref{thm:chiral-qg-equiv}: the $R$-matrix is
given by step~5, the coproduct by steps~4 and~6, and the
$\Ainf$ structure is extracted from
$\Barord(\cW_{1+\infty})$ via the functor $\alpha$
with input the $R$-matrix from step~5.
\end{proof}

\begin{remark}[Structural features of the spectral coproduct]
\label{rem:w-infty-coprod-structure}
The coproduct of
Theorem~\textup{\ref{thm:w-infty-chiral-qg}}
has the following convention-independent features:
\begin{enumerate}[label=\textup{(\arabic*)}]
\item $\Delta_z(J) = J \otimes 1
  + 1 \otimes J$ is cocommutative with no
  $z$-dependence (Heisenberg is abelian,
  confirming class~$G$).
\item $\Delta_z(\psi_2)$ has a noncocommutative
  cross-term $\psi_1 \otimes \psi_1$ at order
  $z^0$ and a spectral tail
  $z\,(1 \otimes \psi_1)$ from the shift
  $u \mapsto u - z$.
  In the $\cW_{1+\infty}$ basis,
  $\Delta_z(T)$ acquires a correction
  $(\Psi - 1)/\Psi \cdot (J \otimes J)$ from the
  Miura inversion $T = \psi_2 - J^2/(2\Psi)$.
\item $\Delta_z(\psi_3)$ has cross-terms
  $\psi_1 \otimes \psi_2 + \psi_2 \otimes \psi_1$,
  a spectral tail through $z^2$, and the $z^1$
  coefficient involves both $\psi_1 \otimes \psi_1$
  and $1 \otimes \psi_2$. In the $\cW_{1+\infty}$
  basis, $\Delta_z(W)$ mixes all three spins with
  coefficients polynomial in $\Psi$.
\item The general formula~\eqref{eq:gl1-coprod-general}
  gives $\Delta_z(\psi_n)$ at every spin in closed form.
  The coproduct on the $\cW_{1+\infty}$ field basis
  $\{J, T, W, W_4, \ldots\}$ is obtained by
  inverting the Miura transform; the computation is
  triangular and algorithmic at each spin.
\end{enumerate}
\end{remark}

\begin{remark}[Effective central charge and intertwining
in the Miura basis]
\label{rem:spin2-ceff-miura}
The coproduct~\eqref{eq:coprod-T} in the
$\cW_{1+\infty}$ field basis has cross-term coefficient
$(\Psi - 1)/\Psi$ (not $1/\Psi$), as derived in step~4.
